\cotefile{63}
\scanpage{73}
\[
\begin{aligned}
 V&=L+L^{-1} & L\ \text{défini par}\quad &T\ \text{\unsure{gpe|gr|sur}}\ U=\{z,\lvert z\rvert=1\}\\
 &&&P\ \text{gpe}\ \mathbb R^{*+}
\end{aligned}
\]
\[
 \varphi:\overline V\xrightarrow{\sim}V\quad\text{involutif}
\]
\[
 \overline L\longrightarrow L+L^{-1}\qquad
 \alpha\in\underbrace{\overline L^{-1}\otimes L}_{\text{défini par }T^{(2)}\ \text{groupe }U},
 \qquad \beta\in\bigl(\underbrace{\overline L L}_{\text{défini par }P^{\unsure{(2)|(0)}}\ \text{gpe}\ \mathbb R^{++}}\bigr)^{-1}
\]

$\boxed{\beta\ne0}$\qquad Alors la \unsure{condition|conditions} \unsure{de} $\alpha$ et $\beta$ \inserted{(par ailleurs arbitraires)} \unsure{détermine|découlent} l'involution car on aura
\[
 V=L+\varphi(L)
\]
\unsure{tt|H} \unsure{élém\textsuperscript{t}|dual} de $V$ s'écrit $z=x+\deleted{\unsure{\varphi(y)}}\bar y$ \quad ($x\in L$, $y\in L$ \unsure{unique|uniques}) et $\bar z=y+\bar x$
\[
 \varphi(\unsure{x|\alpha})=\bar x=\alpha(\bar x)+\beta(\bar x)\qquad
 \deleted{\unsure{x\mapsto\varphi}}
\]
si $y\in L^{-1}$, on aura
\[
\begin{aligned}
 y&=a+\varphi(b) &&a\in L,\ b\in L\\
  &=(a+\alpha(b))+\beta(b)
\end{aligned}
\]
donc $a+\alpha(b)=0$, $\beta(b)=y$, d'où
\[
\begin{aligned}
 b&=\beta^{-1}y,\quad a=-\alpha(b)\\
 &\hphantom{=\beta^{-1}y,\quad a}=-\alpha\beta^{-1}(y)
\end{aligned}
\]
donc
\[
 y=-\alpha\beta^{-1}(y)+\deleted{\unsure{\varphi|\alpha}}\overline{\beta^{-1}(y)},\qquad\text{d'où}
\]
\[
 \begin{aligned}
 &\deleted{\unsure{\varphi(y)=|\varphi(y)}}\\[-2pt]
 \bar y&=\deleted{\unsure{\alpha}}\beta^{-1}(y)-\overline{\alpha\beta(y)}
 =\bigl(\deleted{\unsure{\alpha}}\beta^{-1}(y)\deleted{\unsure{-|+}}\alpha^2\beta^{-1}(y)\bigr)\deleted{\unsure{-|+}}\beta\alpha\beta^{-1}(y):
 \end{aligned}
\]

Sous quelle condition, pour la forme \inserted{alternée} naturelle sur $V$, a-t-on $\langle\bar x,\bar y\rangle=\overline{\langle x,y\rangle}$ ?

Il suffit de l'écrire pour $x\in L$, $y\in L^{-1}$, or alors
\[
 \begin{aligned}
 \langle\bar x,\bar y\rangle
 &=\bigl\langle\underset{L}{\alpha(x)}+\underset{L^{-1}}{\beta(x)},
       \bigl(\deleted{\unsure{\alpha}}\beta^{-1}(y)-\alpha^2\beta^{-1}(y)\bigr)-\beta\alpha\beta^{-1}(y))\\
 &=-\langle\alpha(x),\beta\alpha\beta^{-1}(y)\rangle+
 \langle\beta^{-1}y,-\alpha^{\unsure{2}}\,\beta^{-1}y,\beta(x)\rangle
 \end{aligned}
\]
\uline{On trouve la condition} : $\beta+\bar\beta=0$ \uline{i.e.} $\beta$ \uline{purement imaginaire}
% archive sheet number: 72
